%!TEX root = TheoretischePhysikC.tex
\section{Vorlesung: Drehimpuls}

\subsection{Bahndrehimpuls}
Wir untersuchen nun Quantensysteme, die einen Bahndrehimpuls in dem klassischen Limes aufweisen. Das kanonische Beispiel ist das eines Teilchens, das sich in drei räumlichen Dimensionen bewegen kann. Daher werden wir den Hilbertraum $\mathcal{H} \equiv L^2(\mathbb{R}^2)$ zuordnen um die Zustände zu modellieren. Unser Ziel ist es, die richtigen Quantenobservablen und ihre Eigenschaften zu untersuchen.

Nach unserer Quantisierungstabelle identifizieren wir die folgenden Operatoren als Observablen, die zu der klassischen Drehimpulskomponente gehören.
\begin{definition}
	Die \emph{Bahndrehimpulsoperatoren} lauten $\widehat{\mathbf{L}} \equiv \widehat{\mathbf{r}}\times \widehat{\mathbf{p}}$, wobei 
	\begin{equation}
		\mathbf{r} \equiv \begin{pmatrix}
			\widehat{x} \\
			\widehat{y} \\
			\widehat{z} 
		\end{pmatrix}\quad
		\mathbf{p} \equiv \begin{pmatrix}
			\widehat{p}_x \\
			\widehat{p}_y \\
			\widehat{p}_z 
		\end{pmatrix}\quad\text{und}\quad \widehat{\mathbf{L}} \equiv \begin{pmatrix}
			\widehat{L}_x \\
			\widehat{L}_y \\
			\widehat{L}_z 
		\end{pmatrix} \equiv \begin{pmatrix}
			\widehat{L}_1 \\
			\widehat{L}_2 \\
			\widehat{L}_3 
		\end{pmatrix}.
	\end{equation}
	Die Positionsbasis-Matrixelemente lauten  
	\begin{align}
		\langle x,y,z|\widehat{L}_x|\Psi\rangle &= -i{\hbar} \left(y\frac{\partial}{\partial z} - z\frac{\partial}{\partial y} \right)\Psi(x,y,z) \\
		\langle x,y,z|\widehat{L}_y|\Psi\rangle &=  i{\hbar} \left(x\frac{\partial}{\partial z} - z\frac{\partial}{\partial x} \right)\Psi(x,y,z) \\
		\langle x,y,z|\widehat{L}_z|\Psi\rangle &=  -i{\hbar} \left(x\frac{\partial}{\partial y} - y\frac{\partial}{\partial x} \right)\Psi(x,y,z),
	\end{align}
	wobei $\Psi(x,y,z) = \langle x,y,z|\Psi\rangle$ eine beliebige Wellenfunktion ist. 
\end{definition}

Man kann folgende Kommutatoren zeigen:
\begin{equation}
	\eqbox{[\widehat{L}_1,\widehat{L}_2] = i\hbar \widehat{L}_3, \quad [\widehat{L}_2,\widehat{L}_3] = i\hbar \widehat{L}_1\quad\text{und}\quad [\widehat{L}_3,\widehat{L}_1] = i\hbar \widehat{L}_2.}
\end{equation}
Es folgt auch, dass
\begin{equation}
	\eqbox{[\widehat{\mathbf{L}}^2, \widehat{L}_j]=0, \quad \text{für alle $j=1,2,3$,}}
\end{equation}
wobei
\begin{equation}
	\widehat{\mathbf{L}}^2 \equiv \widehat{L}_1^2+\widehat{L}_2^2+\widehat{L}_3^2.
\end{equation}

Die Kommutatoren werden wie folgt bewiesen: Sei $\Psi(x,y,z)$ eine beliebige Wellenfunktion und betrachten
\begin{equation}
	\begin{split}
		\langle x,y,z|\widehat{L}_1\widehat{L}_2|\Psi\rangle &= {\hbar}^2 \left(y\frac{\partial}{\partial z} - z\frac{\partial}{\partial y} \right)\left(x\frac{\partial}{\partial z} - z\frac{\partial}{\partial x} \right)\Psi(x,y,z) \\
		&= {\hbar}^2 \left(xy\frac{\partial^2}{\partial z^2} - y\frac{\partial}{\partial x} - yz\frac{\partial^2}{\partial x \partial z} - xz\frac{\partial^2}{\partial y\partial z}  + z^2\frac{\partial^2}{\partial x\partial z} \right)\Psi(x,y,z)
	\end{split}
\end{equation}
und 
\begin{equation}
	\begin{split}
		\langle x,y,z|\widehat{L}_2\widehat{L}_1|\Psi\rangle &= {\hbar}^2 \left(x\frac{\partial}{\partial z} - z\frac{\partial}{\partial x} \right)\left(y\frac{\partial}{\partial z} - z\frac{\partial}{\partial y} \right)\Psi(x,y,z) \\
		&= {\hbar}^2 \left(xy\frac{\partial^2}{\partial z^2} - x\frac{\partial}{\partial y} - xz\frac{\partial^2}{\partial y \partial z} - yz\frac{\partial^2}{\partial x\partial z}  + z^2\frac{\partial^2}{\partial x\partial z} \right)\Psi(x,y,z).
	\end{split}
\end{equation}
Daher berechnen wir den Kommutator
\begin{equation}
	\langle x,y,z|\widehat{L}_1\widehat{L}_2 - \widehat{L}_2\widehat{L}_1|\Psi\rangle = {\hbar}^2 \left(- y\frac{\partial}{\partial x} + x\frac{\partial}{\partial y} \right)\Psi(x,y,z) = i\hbar \langle x,y,z|\widehat{L}_3|\Psi\rangle.
\end{equation}

Man kann beweisen, dass $\widehat{\mathbf{L}}^2$ und $\widehat{L}_3$ besitzen gemeinsame Eigenzustände. Diese Eigenstände werden $|\ell,m\rangle$ bezeichnet und erfüllen die folgende Eigenwertgleichungen 
\begin{equation}
	\eqbox{
	\begin{aligned}
		\widehat{\mathbf{L}}^2|\ell,m\rangle &= \hbar^2\ell(\ell+1)|\ell,m\rangle, \\
		\widehat{L}_3|\ell,m\rangle &= \hbar m|\ell,m\rangle,
	\end{aligned}
	}
\end{equation}
wobei $\ell = 0, \frac12, 1, \frac32, \ldots$ und $m = -\ell, -\ell + 1, \ldots, \ell$. Hierbei bezeichnet $\ell$ die \emph{Drehimpulsquantenzahl} und $m$ die \emph{magnetische Quantenzahl}.

Nun werden wir die Orstdarstellung der Eigenfunktionen herleiten. Mit Hilfe von den Kugelkoordinaten $r$, $\theta$, $\phi$ erhalten wir (Übung)
\begin{equation}
	\begin{split}
		\langle r,\theta,\phi|\widehat{L}_3|\Psi\rangle &= -i\hbar \frac{\partial}{\partial \phi} \Psi(r,\theta,\phi), \\
		\langle r,\theta,\phi|\widehat{\mathbf{L}}^2|\Psi\rangle &= -\hbar^2\frac{1}{\sin^2(\theta)}\left(\sin(\theta)\frac{\partial}{\partial \theta}\left(\sin(\theta)\frac{\partial}{\partial \theta}\right) + \frac{\partial^2}{\partial \phi^2}\right) \Psi(r,\theta,\phi),
	\end{split}
\end{equation}
wobei $|\Psi\rangle \in L^2(\mathbb{R}^3)$ einen beliebigen reinen Zustand ist. 

Die Eigenfunktionen von $\widehat{L}_3$ sind durch $e^{im\phi}$ gegeben:
\begin{equation}
	-i\hbar \frac{\partial}{\partial \phi}e^{im\phi} = \hbar m e^{im\phi}.
\end{equation}
Die Eindeutigkeit der Wellenfunktion $\Psi(r,\theta,\phi)$ unter die Verschiebung $\phi\mapsto \phi+2\pi$ fordert $m\in \mathbb{Z}$, sodass nur $\ell = 0,1,2, \ldots$ auftritt. (Die Eigenvektoren für Halbzahlige $\ell$ wird später mit eine andere Methode aufgebaut.) Die gemeinsamen Eigenfunktionen von $\widehat{\mathbf{L}}^2$ und $\widehat{L}_3$ lauten
\begin{equation}
	\eqbox{\text{Y}_\ell^m (\theta,\phi) \equiv (-1)^m\left(\frac{2\ell + 1}{4\pi}\frac{(\ell-m)!}{(\ell+m)!}\right)^{\frac12} \text{P}_\ell^m(\cos(\theta))e^{im\phi},}
\end{equation}
wobei für $m\ge 0$
\begin{equation}
	\eqbox{\text{P}^m_\ell(x) = (1-x^2)^{\frac{m}{2}}\frac{d^m}{dx^m}\text{P}_\ell(x)}
\end{equation}
die \emph{assoziierten Legendrefunktionen} die mit Hilfe von den \emph{Legendre-Polynomen}
\begin{equation}
	\eqbox{\text{P}_\ell(x)=\frac{1}{2^\ell \ell!} \frac{d^\ell}{dx^\ell}(x^2-1)^\ell}
\end{equation}
definiert sind. Die Funktionen $\text{Y}_\ell^m$ sind als die \emph{Kugelflächenfunktionen} bekannt. 

Die ersten drei Kugelflächenfunktionen lauten
\begin{align}
	\text{Y}^0_0(\theta,\phi) &= \frac{1}{\sqrt{4\pi}}, \\
	\text{Y}^0_1(\theta,\phi) &= \sqrt{\frac{3}{4\pi}}\cos(\theta), \\
	\text{Y}^1_1(\theta,\phi) &= -\sqrt{\frac{3}{8\pi}}\sin(\theta)e^{i\phi}.
\end{align}